\IfFileExists{stacks-project.cls}{%
\documentclass{stacks-project}
}{%
\documentclass{amsart}
}

% For dealing with references we use the comment environment
\usepackage{verbatim}
\newenvironment{reference}{\comment}{\endcomment}
%\newenvironment{reference}{}{}
\newenvironment{slogan}{\comment}{\endcomment}
\newenvironment{history}{\comment}{\endcomment}

% For commutative diagrams we use Xy-pic
\usepackage[all]{xy}

% We use 2cell for 2-commutative diagrams.
\xyoption{2cell}
\UseAllTwocells

% We use multicol for the list of chapters between chapters
\usepackage{multicol}

% This is generally recommended for better output
\usepackage{lmodern}
\usepackage[T1]{fontenc}

% For cross-file-references
\usepackage{xr-hyper}

% Package for hypertext links:
\usepackage{hyperref}

% For any local file, say "hello.tex" you want to link to please
% use \externaldocument[hello-]{hello}
\externaldocument[introduction-]{introduction}
\externaldocument[conventions-]{conventions}
\externaldocument[sets-]{sets}
\externaldocument[categories-]{categories}
\externaldocument[topology-]{topology}
\externaldocument[sheaves-]{sheaves}
\externaldocument[sites-]{sites}
\externaldocument[stacks-]{stacks}
\externaldocument[fields-]{fields}
\externaldocument[algebra-]{algebra}
\externaldocument[brauer-]{brauer}
\externaldocument[homology-]{homology}
\externaldocument[derived-]{derived}
\externaldocument[simplicial-]{simplicial}
\externaldocument[more-algebra-]{more-algebra}
\externaldocument[smoothing-]{smoothing}
\externaldocument[modules-]{modules}
\externaldocument[sites-modules-]{sites-modules}
\externaldocument[injectives-]{injectives}
\externaldocument[cohomology-]{cohomology}
\externaldocument[sites-cohomology-]{sites-cohomology}
\externaldocument[dga-]{dga}
\externaldocument[dpa-]{dpa}
\externaldocument[sdga-]{sdga}
\externaldocument[hypercovering-]{hypercovering}
\externaldocument[schemes-]{schemes}
\externaldocument[constructions-]{constructions}
\externaldocument[properties-]{properties}
\externaldocument[morphisms-]{morphisms}
\externaldocument[coherent-]{coherent}
\externaldocument[divisors-]{divisors}
\externaldocument[limits-]{limits}
\externaldocument[varieties-]{varieties}
\externaldocument[topologies-]{topologies}
\externaldocument[descent-]{descent}
\externaldocument[perfect-]{perfect}
\externaldocument[more-morphisms-]{more-morphisms}
\externaldocument[flat-]{flat}
\externaldocument[groupoids-]{groupoids}
\externaldocument[more-groupoids-]{more-groupoids}
\externaldocument[etale-]{etale}
\externaldocument[chow-]{chow}
\externaldocument[intersection-]{intersection}
\externaldocument[pic-]{pic}
\externaldocument[weil-]{weil}
\externaldocument[adequate-]{adequate}
\externaldocument[dualizing-]{dualizing}
\externaldocument[duality-]{duality}
\externaldocument[discriminant-]{discriminant}
\externaldocument[derham-]{derham}
\externaldocument[local-cohomology-]{local-cohomology}
\externaldocument[algebraization-]{algebraization}
\externaldocument[curves-]{curves}
\externaldocument[resolve-]{resolve}
\externaldocument[models-]{models}
\externaldocument[functors-]{functors}
\externaldocument[equiv-]{equiv}
\externaldocument[pione-]{pione}
\externaldocument[etale-cohomology-]{etale-cohomology}
\externaldocument[proetale-]{proetale}
\externaldocument[relative-cycles-]{relative-cycles}
\externaldocument[more-etale-]{more-etale}
\externaldocument[trace-]{trace}
\externaldocument[crystalline-]{crystalline}
\externaldocument[spaces-]{spaces}
\externaldocument[spaces-properties-]{spaces-properties}
\externaldocument[spaces-morphisms-]{spaces-morphisms}
\externaldocument[decent-spaces-]{decent-spaces}
\externaldocument[spaces-cohomology-]{spaces-cohomology}
\externaldocument[spaces-limits-]{spaces-limits}
\externaldocument[spaces-divisors-]{spaces-divisors}
\externaldocument[spaces-over-fields-]{spaces-over-fields}
\externaldocument[spaces-topologies-]{spaces-topologies}
\externaldocument[spaces-descent-]{spaces-descent}
\externaldocument[spaces-perfect-]{spaces-perfect}
\externaldocument[spaces-more-morphisms-]{spaces-more-morphisms}
\externaldocument[spaces-flat-]{spaces-flat}
\externaldocument[spaces-groupoids-]{spaces-groupoids}
\externaldocument[spaces-more-groupoids-]{spaces-more-groupoids}
\externaldocument[bootstrap-]{bootstrap}
\externaldocument[spaces-pushouts-]{spaces-pushouts}
\externaldocument[spaces-chow-]{spaces-chow}
\externaldocument[groupoids-quotients-]{groupoids-quotients}
\externaldocument[spaces-more-cohomology-]{spaces-more-cohomology}
\externaldocument[spaces-simplicial-]{spaces-simplicial}
\externaldocument[spaces-duality-]{spaces-duality}
\externaldocument[formal-spaces-]{formal-spaces}
\externaldocument[restricted-]{restricted}
\externaldocument[spaces-resolve-]{spaces-resolve}
\externaldocument[formal-defos-]{formal-defos}
\externaldocument[defos-]{defos}
\externaldocument[cotangent-]{cotangent}
\externaldocument[examples-defos-]{examples-defos}
\externaldocument[algebraic-]{algebraic}
\externaldocument[examples-stacks-]{examples-stacks}
\externaldocument[stacks-sheaves-]{stacks-sheaves}
\externaldocument[criteria-]{criteria}
\externaldocument[artin-]{artin}
\externaldocument[quot-]{quot}
\externaldocument[stacks-properties-]{stacks-properties}
\externaldocument[stacks-morphisms-]{stacks-morphisms}
\externaldocument[stacks-limits-]{stacks-limits}
\externaldocument[stacks-cohomology-]{stacks-cohomology}
\externaldocument[stacks-perfect-]{stacks-perfect}
\externaldocument[stacks-introduction-]{stacks-introduction}
\externaldocument[stacks-more-morphisms-]{stacks-more-morphisms}
\externaldocument[stacks-geometry-]{stacks-geometry}
\externaldocument[moduli-]{moduli}
\externaldocument[moduli-curves-]{moduli-curves}
\externaldocument[examples-]{examples}
\externaldocument[exercises-]{exercises}
\externaldocument[guide-]{guide}
\externaldocument[desirables-]{desirables}
\externaldocument[coding-]{coding}
\externaldocument[obsolete-]{obsolete}
\externaldocument[fdl-]{fdl}
\externaldocument[index-]{index}

% Theorem environments.
%
\theoremstyle{plain}
\newtheorem{theorem}[subsection]{Theorem}
\newtheorem{proposition}[subsection]{Proposition}
\newtheorem{lemma}[subsection]{Lemma}

\theoremstyle{definition}
\newtheorem{definition}[subsection]{Definition}
\newtheorem{example}[subsection]{Example}
\newtheorem{exercise}[subsection]{Exercise}
\newtheorem{situation}[subsection]{Situation}

\theoremstyle{remark}
\newtheorem{remark}[subsection]{Remark}
\newtheorem{remarks}[subsection]{Remarks}

\numberwithin{equation}{subsection}

% Macros
%
\def\lim{\mathop{\mathrm{lim}}\nolimits}
\def\colim{\mathop{\mathrm{colim}}\nolimits}
\def\Spec{\mathop{\mathrm{Spec}}}
\def\Hom{\mathop{\mathrm{Hom}}\nolimits}
\def\Ext{\mathop{\mathrm{Ext}}\nolimits}
\def\SheafHom{\mathop{\mathcal{H}\!\mathit{om}}\nolimits}
\def\SheafExt{\mathop{\mathcal{E}\!\mathit{xt}}\nolimits}
\def\Sch{\mathit{Sch}}
\def\Mor{\mathop{\mathrm{Mor}}\nolimits}
\def\Ob{\mathop{\mathrm{Ob}}\nolimits}
\def\Sh{\mathop{\mathit{Sh}}\nolimits}
\def\NL{\mathop{N\!L}\nolimits}
\def\CH{\mathop{\mathrm{CH}}\nolimits}
\def\proetale{{pro\text{-}\acute{e}tale}}
\def\etale{{\acute{e}tale}}
\def\QCoh{\mathit{QCoh}}
\def\Ker{\mathop{\mathrm{Ker}}}
\def\Im{\mathop{\mathrm{Im}}}
\def\Coker{\mathop{\mathrm{Coker}}}
\def\Coim{\mathop{\mathrm{Coim}}}

% Boxtimes
%
\DeclareMathSymbol{\boxtimes}{\mathbin}{AMSa}{"02}

%
% Macros for moduli stacks/spaces
%
\def\QCohstack{\mathcal{QC}\!\mathit{oh}}
\def\Cohstack{\mathcal{C}\!\mathit{oh}}
\def\Spacesstack{\mathcal{S}\!\mathit{paces}}
\def\Quotfunctor{\mathrm{Quot}}
\def\Hilbfunctor{\mathrm{Hilb}}
\def\Curvesstack{\mathcal{C}\!\mathit{urves}}
\def\Polarizedstack{\mathcal{P}\!\mathit{olarized}}
\def\Complexesstack{\mathcal{C}\!\mathit{omplexes}}
% \Pic is the operator that assigns to X its picard group, usage \Pic(X)
% \Picardstack_{X/B} denotes the Picard stack of X over B
% \Picardfunctor_{X/B} denotes the Picard functor of X over B
\def\Pic{\mathop{\mathrm{Pic}}\nolimits}
\def\Picardstack{\mathcal{P}\!\mathit{ic}}
\def\Picardfunctor{\mathrm{Pic}}
\def\Deformationcategory{\mathcal{D}\!\mathit{ef}}
\usepackage{cleveref}
\crefname{originallemma}{original lemma}{original lemmas}
\Crefname{originallemma}{Original lemma}{Original lemmas}
\makeatletter
\newlabel{sites-cohomology-remark-base-change}{{07A7}{source}{}{sourcepointer-1}{}}
\newlabel{sites-cohomology-remark-base-change@cref}{{[originallemma][0][]07A7}{[1][0][]source}{}{}{}}
\makeatother
\crefname{originallemma}{original lemma}{original lemmas}
\Crefname{originallemma}{Original lemma}{Original lemmas}
\makeatletter
\newlabel{spaces-cohomology-lemma-finite-higher-direct-image-zero}{{0A4K}{source}{}{sourcepointer-2}{}}
\newlabel{spaces-cohomology-lemma-finite-higher-direct-image-zero@cref}{{[originallemma][0][]0A4K}{[1][0][]source}{}{}{}}
\makeatother
\crefname{originallemma}{original lemma}{original lemmas}
\Crefname{originallemma}{Original lemma}{Original lemmas}
\makeatletter
\newlabel{spaces-topologies-lemma-characterize-sheaf}{{0DFM}{source}{}{sourcepointer-3}{}}
\newlabel{spaces-topologies-lemma-characterize-sheaf@cref}{{[originallemma][0][]0DFM}{[1][0][]source}{}{}{}}
\makeatother
\crefname{originallemma}{original lemma}{original lemmas}
\Crefname{originallemma}{Original lemma}{Original lemmas}
\makeatletter
\newlabel{spaces-topologies-lemma-zariski-etale-smooth-syntomic-fppf-ph}{{0DFH}{source}{}{sourcepointer-4}{}}
\newlabel{spaces-topologies-lemma-zariski-etale-smooth-syntomic-fppf-ph@cref}{{[originallemma][0][]0DFH}{[1][0][]source}{}{}{}}
\makeatother
\crefname{originallemma}{original lemma}{original lemmas}
\Crefname{originallemma}{Original lemma}{Original lemmas}
\makeatletter
\newlabel{etale-cohomology-lemma-describe-pullback}{{09XN}{source}{}{sourcepointer-5}{}}
\newlabel{etale-cohomology-lemma-describe-pullback@cref}{{[originallemma][0][]09XN}{[1][0][]source}{}{}{}}
\makeatother
\crefname{originallemma}{original lemma}{original lemmas}
\Crefname{originallemma}{Original lemma}{Original lemmas}
\makeatletter
\newlabel{spaces-cohomology-lemma-weak-chow}{{089J}{source}{}{sourcepointer-6}{}}
\newlabel{spaces-cohomology-lemma-weak-chow@cref}{{[originallemma][0][]089J}{[1][0][]source}{}{}{}}
\makeatother
\crefname{originallemma}{original lemma}{original lemmas}
\Crefname{originallemma}{Original lemma}{Original lemmas}
\makeatletter
\newlabel{etale-cohomology-lemma-proper-base-change-f-star}{{0A3U}{source}{}{sourcepointer-7}{}}
\newlabel{etale-cohomology-lemma-proper-base-change-f-star@cref}{{[originallemma][0][]0A3U}{[1][0][]source}{}{}{}}
\makeatother
\crefname{originallemma}{original lemma}{original lemmas}
\Crefname{originallemma}{Original lemma}{Original lemmas}
\makeatletter
\newlabel{topologies-lemma-morphism-big-small-cartesian-diagram-etale}{{0DDA}{source}{}{sourcepointer-8}{}}
\newlabel{topologies-lemma-morphism-big-small-cartesian-diagram-etale@cref}{{[originallemma][0][]0DDA}{[1][0][]source}{}{}{}}
\makeatother
\crefname{originallemma}{original lemma}{original lemmas}
\Crefname{originallemma}{Original lemma}{Original lemmas}
\makeatletter
\newlabel{sites-cohomology-lemma-localize-cartesian-square}{{0EZ0}{source}{}{sourcepointer-9}{}}
\newlabel{sites-cohomology-lemma-localize-cartesian-square@cref}{{[originallemma][0][]0EZ0}{[1][0][]source}{}{}{}}
\makeatother
\crefname{originallemma}{original lemma}{original lemmas}
\Crefname{originallemma}{Original lemma}{Original lemmas}
\makeatletter
\newlabel{sites-cohomology-lemma-localize-cartesian-square-modules}{{0FN5}{source}{}{sourcepointer-10}{}}
\newlabel{sites-cohomology-lemma-localize-cartesian-square-modules@cref}{{[originallemma][0][]0FN5}{[1][0][]source}{}{}{}}
\makeatother
\crefname{originallemma}{original lemma}{original lemmas}
\Crefname{originallemma}{Original lemma}{Original lemmas}
\makeatletter
\newlabel{sites-cohomology-lemma-derived-pushforward-composition}{{0D6E}{source}{}{sourcepointer-11}{}}
\newlabel{sites-cohomology-lemma-derived-pushforward-composition@cref}{{[originallemma][0][]0D6E}{[1][0][]source}{}{}{}}
\makeatother
\crefname{originallemma}{original lemma}{original lemmas}
\Crefname{originallemma}{Original lemma}{Original lemmas}
\makeatletter
\newlabel{derived-lemma-two-ss-complex-functor}{{015J}{source}{}{sourcepointer-12}{}}
\newlabel{derived-lemma-two-ss-complex-functor@cref}{{[originallemma][0][]015J}{[1][0][]source}{}{}{}}
\makeatother
\crefname{originallemma}{original lemma}{original lemmas}
\Crefname{originallemma}{Original lemma}{Original lemmas}
\makeatletter
\newlabel{spaces-properties-lemma-stalk-pullback}{{04K2}{source}{}{sourcepointer-13}{}}
\newlabel{spaces-properties-lemma-stalk-pullback@cref}{{[originallemma][0][]04K2}{[1][0][]source}{}{}{}}
\makeatother
\crefname{originallemma}{original lemma}{original lemmas}
\Crefname{originallemma}{Original lemma}{Original lemmas}
\makeatletter
\newlabel{topologies-lemma-surjective-proper-ph}{{0DES}{source}{}{sourcepointer-14}{}}
\newlabel{topologies-lemma-surjective-proper-ph@cref}{{[originallemma][0][]0DES}{[1][0][]source}{}{}{}}
\makeatother
\crefname{originallemma}{original lemma}{original lemmas}
\Crefname{originallemma}{Original lemma}{Original lemmas}
\makeatletter
\newlabel{etale-cohomology-lemma-describe-pullback-pi-ph}{{0DDW}{source}{}{sourcepointer-15}{}}
\newlabel{etale-cohomology-lemma-describe-pullback-pi-ph@cref}{{[originallemma][0][]0DDW}{[1][0][]source}{}{}{}}
\makeatother
\crefname{originallemma}{original lemma}{original lemmas}
\Crefname{originallemma}{Original lemma}{Original lemmas}
\makeatletter
\newlabel{homology-lemma-adjoint-preserve-injectives}{{015Z}{source}{}{sourcepointer-16}{}}
\newlabel{homology-lemma-adjoint-preserve-injectives@cref}{{[originallemma][0][]015Z}{[1][0][]source}{}{}{}}
\makeatother
\crefname{originallemma}{original lemma}{original lemmas}
\Crefname{originallemma}{Original lemma}{Original lemmas}
\makeatletter
\newlabel{derived-lemma-adjoint-preserve-K-injectives}{{08BJ}{source}{}{sourcepointer-17}{}}
\newlabel{derived-lemma-adjoint-preserve-K-injectives@cref}{{[originallemma][0][]08BJ}{[1][0][]source}{}{}{}}
\makeatother
\crefname{originallemma}{original lemma}{original lemmas}
\Crefname{originallemma}{Original lemma}{Original lemmas}
\makeatletter
\newlabel{topologies-lemma-put-in-T-etale}{{021F}{source}{}{sourcepointer-18}{}}
\newlabel{topologies-lemma-put-in-T-etale@cref}{{[originallemma][0][]021F}{[1][0][]source}{}{}{}}
\makeatother
\crefname{originallemma}{original lemma}{original lemmas}
\Crefname{originallemma}{Original lemma}{Original lemmas}
\makeatletter
\newlabel{sites-modules-lemma-g-shriek-adjoint}{{04BG}{source}{}{sourcepointer-19}{}}
\newlabel{sites-modules-lemma-g-shriek-adjoint@cref}{{[originallemma][0][]04BG}{[1][0][]source}{}{}{}}
\makeatother
\crefname{originallemma}{original lemma}{original lemmas}
\Crefname{originallemma}{Original lemma}{Original lemmas}
\makeatletter
\newlabel{sites-modules-lemma-exactness-lower-shriek}{{04BH}{source}{}{sourcepointer-20}{}}
\newlabel{sites-modules-lemma-exactness-lower-shriek@cref}{{[originallemma][0][]04BH}{[1][0][]source}{}{}{}}
\makeatother
\crefname{originallemma}{original lemma}{original lemmas}
\Crefname{originallemma}{Original lemma}{Original lemmas}
\makeatletter
\newlabel{topologies-lemma-morphism-big-small-etale}{{021I}{source}{}{sourcepointer-21}{}}
\newlabel{topologies-lemma-morphism-big-small-etale@cref}{{[originallemma][0][]021I}{[1][0][]source}{}{}{}}
\makeatother
\crefname{originallemma}{original lemma}{original lemmas}
\Crefname{originallemma}{Original lemma}{Original lemmas}
\makeatletter
\newlabel{sites-cohomology-lemma-cohomology-of-open}{{03F3}{source}{}{sourcepointer-22}{}}
\newlabel{sites-cohomology-lemma-cohomology-of-open@cref}{{[originallemma][0][]03F3}{[1][0][]source}{}{}{}}
\makeatother
\crefname{originallemma}{original lemma}{original lemmas}
\Crefname{originallemma}{Original lemma}{Original lemmas}
\makeatletter
\newlabel{sites-cohomology-lemma-restrict-K-injective-to-open}{{08FI}{source}{}{sourcepointer-23}{}}
\newlabel{sites-cohomology-lemma-restrict-K-injective-to-open@cref}{{[originallemma][0][]08FI}{[1][0][]source}{}{}{}}
\makeatother
\crefname{originallemma}{original lemma}{original lemmas}
\Crefname{originallemma}{Original lemma}{Original lemmas}
\makeatletter
\newlabel{etale-cohomology-lemma-compare-injectives}{{0758}{source}{}{sourcepointer-24}{}}
\newlabel{etale-cohomology-lemma-compare-injectives@cref}{{[originallemma][0][]0758}{[1][0][]source}{}{}{}}
\makeatother
\crefname{originallemma}{original lemma}{original lemmas}
\Crefname{originallemma}{Original lemma}{Original lemmas}
\makeatletter
\newlabel{etale-cohomology-section-proper-base-change}{{095S}{source}{}{sourcepointer-25}{}}
\newlabel{etale-cohomology-section-proper-base-change@cref}{{[originallemma][0][]095S}{[1][0][]source}{}{}{}}
\makeatother
\crefname{originallemma}{original lemma}{original lemmas}
\Crefname{originallemma}{Original lemma}{Original lemmas}
\makeatletter
\newlabel{spaces-cohomology-lemma-colimit-cohomology}{{07U6}{source}{}{sourcepointer-26}{}}
\newlabel{spaces-cohomology-lemma-colimit-cohomology@cref}{{[originallemma][0][]07U6}{[1][0][]source}{}{}{}}
\makeatother
\crefname{originallemma}{original lemma}{original lemmas}
\Crefname{originallemma}{Original lemma}{Original lemmas}
\makeatletter
\newlabel{derived-lemma-leray-acyclicity}{{015E}{source}{}{sourcepointer-27}{}}
\newlabel{derived-lemma-leray-acyclicity@cref}{{[originallemma][0][]015E}{[1][0][]source}{}{}{}}
\makeatother
\title{Proper comparison of higher direct images: pi\_X\textasciicircum{}\{-1\} R f\_\{small,*\} = R f\_\{big,*\} pi\_Y\textasciicircum{}\{-1\}}
\author{The Stacks Project Authors}
\date{stacks-project:spaces-more-cohomology.tex:local-snapshot-20261012; GNU FDL 1.2 or later}
\begin{document}
\maketitle
\noindent\textit{Editorial note.} The counted statement is the author's own lemma environment labelled lemma-compare-higher-direct-image-proper (source lines 1014-1027); the proof environment beginning on the line after the following blank line (source line 1029) is its own proof and closes at source line 1099.\par\medskip
\noindent\textit{Editorial note.} The proof invokes four results of the same chapter by label, all retained as context with the generator's transitive expansion; two cross-chapter results (Cohomology on Sites, Cohomology of Spaces) are declared as tagged source pointers.\par\medskip
\noindent\textit{Editorial note.} The author's own text refers to the fibre product diagram and to the objects g : X' -> X of the small etale site; both are reproduced verbatim, and nothing is added to the argument.\par\medskip
\noindent\textit{Editorial note.} The retained preamble is the chapter's own (sources/stacks-spaces-more-cohomology/preamble.tex, taken with its real line count), a verbatim excerpt of the native Xy-pic setup and theorem declarations; no macro was invented or redefined.\par\medskip
\noindent\textit{Editorial scope.} One original result of the Stacks Project's chapter Cohomology of Algebraic Spaces, counted once: the comparison of the higher direct images of a proper morphism against the big and small etale topologies, in its underived, derived and finite forms.\par\medskip
\noindent\textit{Allowed standard prerequisites (editorial).}\begin{itemize}
\item The description of pi\textasciicircum{}\{-1\}. pi\_X\textasciicircum{}\{-1\}F is computed on an object X' of the big etale site by Gamma(X'\_etale, (g\_\{small\})\textasciicircum{}\{-1\}F), and similarly for pi\_Y (the author's lemma-describe-pullback, retained as context).. It gives the second equality of part (1) and the second displayed equality at source lines 1046 and 1052.
\item Proper base change for sheaves and complexes. The proper base change isomorphism for sheaves of sets and the base change theorem for complexes with torsion cohomology sheaves (the author's lemma-proper-base-change-f-star and lemma-proper-base-change, retained as context).. They finish parts (1) and (2) at source lines 1053-1055 and 1086-1088.
\item The comparison of underived and derived pushforwards. The comparison between f\_\{small,*\} and f\_\{big,*\} and its parts (1) and (3) for the base changed morphism (the author's lemma-compare-higher-direct-image, retained as context).. It rewrites the pullback of the base change map at source lines 1083-1085 and gives the exactness of f\_\{big,*\} at source lines 1095-1097.
\item Exactness of the direct image for a finite morphism. For f finite the functor f\_\{small,*\} is exact, hence its higher direct images vanish (Cohomology of Spaces, declared as a tagged source pointer).. It is the first step of part (3) at source lines 1091-1094.
\item The canonical base change map. The canonical map pi\_X\textasciicircum{}\{-1\}Rf\_\{small,*\}K -> Rf\_\{big,*\}pi\_Y\textasciicircum{}\{-1\}K whose being an isomorphism is the content of part (2) (Cohomology on Sites, declared as a tagged source pointer).. It is the map studied in part (2) at source lines 1063-1066.
\end{itemize}\medskip
\ifdefined\hypertarget\hypertarget{sourcepointer-1}{}\fi
\noindent\textit{Source reference.} Original lemma 07A7: sources/stacks-chow/tags.txt, source lines 9297--9297. Cross-chapter pointer whose target carries the canonical tag 07A7 in the project tag list; declared so the reference resolves inside the assembled package, while every step the counted argument uses remains inside the retained spans.\par\medskip
\ifdefined\hypertarget\hypertarget{sourcepointer-2}{}\fi
\noindent\textit{Source reference.} Original lemma 0A4K: sources/stacks-chow/tags.txt, source lines 12778--12778. Cross-chapter pointer whose target carries the canonical tag 0A4K in the project tag list; declared so the reference resolves inside the assembled package, while every step the counted argument uses remains inside the retained spans.\par\medskip
\ifdefined\hypertarget\hypertarget{sourcepointer-3}{}\fi
\noindent\textit{Source reference.} Original lemma 0DFM: sources/stacks-chow/tags.txt, source lines 16842--16842. Cross-chapter pointer whose target carries the canonical tag 0DFM in the project tag list; declared so the reference resolves inside the assembled package, while every step the counted argument uses remains inside the retained spans.\par\medskip
\ifdefined\hypertarget\hypertarget{sourcepointer-4}{}\fi
\noindent\textit{Source reference.} Original lemma 0DFH: sources/stacks-chow/tags.txt, source lines 16837--16837. Cross-chapter pointer whose target carries the canonical tag 0DFH in the project tag list; declared so the reference resolves inside the assembled package, while every step the counted argument uses remains inside the retained spans.\par\medskip
\ifdefined\hypertarget\hypertarget{sourcepointer-5}{}\fi
\noindent\textit{Source reference.} Original lemma 09XN: sources/stacks-chow/tags.txt, source lines 12536--12536. Cross-chapter pointer whose target carries the canonical tag 09XN in the project tag list; declared so the reference resolves inside the assembled package, while every step the counted argument uses remains inside the retained spans.\par\medskip
\ifdefined\hypertarget\hypertarget{sourcepointer-6}{}\fi
\noindent\textit{Source reference.} Original lemma 089J: sources/stacks-chow/tags.txt, source lines 10499--10499. Cross-chapter pointer whose target carries the canonical tag 089J in the project tag list; declared so the reference resolves inside the assembled package, while every step the counted argument uses remains inside the retained spans.\par\medskip
\ifdefined\hypertarget\hypertarget{sourcepointer-7}{}\fi
\noindent\textit{Source reference.} Original lemma 0A3U: sources/stacks-chow/tags.txt, source lines 12752--12752. Cross-chapter pointer whose target carries the canonical tag 0A3U in the project tag list; declared so the reference resolves inside the assembled package, while every step the counted argument uses remains inside the retained spans.\par\medskip
\ifdefined\hypertarget\hypertarget{sourcepointer-8}{}\fi
\noindent\textit{Source reference.} Original lemma 0DDA: sources/stacks-chow/tags.txt, source lines 16760--16760. Cross-chapter pointer whose target carries the canonical tag 0DDA in the project tag list; declared so the reference resolves inside the assembled package, while every step the counted argument uses remains inside the retained spans.\par\medskip
\ifdefined\hypertarget\hypertarget{sourcepointer-9}{}\fi
\noindent\textit{Source reference.} Original lemma 0EZ0: sources/stacks-chow/tags.txt, source lines 18710--18710. Cross-chapter pointer whose target carries the canonical tag 0EZ0 in the project tag list; declared so the reference resolves inside the assembled package, while every step the counted argument uses remains inside the retained spans.\par\medskip
\ifdefined\hypertarget\hypertarget{sourcepointer-10}{}\fi
\noindent\textit{Source reference.} Original lemma 0FN5: sources/stacks-chow/tags.txt, source lines 19555--19555. Cross-chapter pointer whose target carries the canonical tag 0FN5 in the project tag list; declared so the reference resolves inside the assembled package, while every step the counted argument uses remains inside the retained spans.\par\medskip
\ifdefined\hypertarget\hypertarget{sourcepointer-11}{}\fi
\noindent\textit{Source reference.} Original lemma 0D6E: sources/stacks-chow/tags.txt, source lines 16519--16519. Cross-chapter pointer whose target carries the canonical tag 0D6E in the project tag list; declared so the reference resolves inside the assembled package, while every step the counted argument uses remains inside the retained spans.\par\medskip
\ifdefined\hypertarget\hypertarget{sourcepointer-12}{}\fi
\noindent\textit{Source reference.} Original lemma 015J: sources/stacks-chow/tags.txt, source lines 1510--1510. Cross-chapter pointer whose target carries the canonical tag 015J in the project tag list; declared so the reference resolves inside the assembled package, while every step the counted argument uses remains inside the retained spans.\par\medskip
\ifdefined\hypertarget\hypertarget{sourcepointer-13}{}\fi
\noindent\textit{Source reference.} Original lemma 04K2: sources/stacks-chow/tags.txt, source lines 5958--5958. Cross-chapter pointer whose target carries the canonical tag 04K2 in the project tag list; declared so the reference resolves inside the assembled package, while every step the counted argument uses remains inside the retained spans.\par\medskip
\ifdefined\hypertarget\hypertarget{sourcepointer-14}{}\fi
\noindent\textit{Source reference.} Original lemma 0DES: sources/stacks-chow/tags.txt, source lines 16812--16812. Cross-chapter pointer whose target carries the canonical tag 0DES in the project tag list; declared so the reference resolves inside the assembled package, while every step the counted argument uses remains inside the retained spans.\par\medskip
\ifdefined\hypertarget\hypertarget{sourcepointer-15}{}\fi
\noindent\textit{Source reference.} Original lemma 0DDW: sources/stacks-chow/tags.txt, source lines 16781--16781. Cross-chapter pointer whose target carries the canonical tag 0DDW in the project tag list; declared so the reference resolves inside the assembled package, while every step the counted argument uses remains inside the retained spans.\par\medskip
\ifdefined\hypertarget\hypertarget{sourcepointer-16}{}\fi
\noindent\textit{Source reference.} Original lemma 015Z: sources/stacks-chow/tags.txt, source lines 1526--1526. Cross-chapter pointer whose target carries the canonical tag 015Z in the project tag list; declared so the reference resolves inside the assembled package, while every step the counted argument uses remains inside the retained spans.\par\medskip
\ifdefined\hypertarget\hypertarget{sourcepointer-17}{}\fi
\noindent\textit{Source reference.} Original lemma 08BJ: sources/stacks-chow/tags.txt, source lines 10569--10569. Cross-chapter pointer whose target carries the canonical tag 08BJ in the project tag list; declared so the reference resolves inside the assembled package, while every step the counted argument uses remains inside the retained spans.\par\medskip
\ifdefined\hypertarget\hypertarget{sourcepointer-18}{}\fi
\noindent\textit{Source reference.} Original lemma 021F: sources/stacks-chow/tags.txt, source lines 2668--2668. Cross-chapter pointer whose target carries the canonical tag 021F in the project tag list; declared so the reference resolves inside the assembled package, while every step the counted argument uses remains inside the retained spans.\par\medskip
\ifdefined\hypertarget\hypertarget{sourcepointer-19}{}\fi
\noindent\textit{Source reference.} Original lemma 04BG: sources/stacks-chow/tags.txt, source lines 5651--5651. Cross-chapter pointer whose target carries the canonical tag 04BG in the project tag list; declared so the reference resolves inside the assembled package, while every step the counted argument uses remains inside the retained spans.\par\medskip
\ifdefined\hypertarget\hypertarget{sourcepointer-20}{}\fi
\noindent\textit{Source reference.} Original lemma 04BH: sources/stacks-chow/tags.txt, source lines 5652--5652. Cross-chapter pointer whose target carries the canonical tag 04BH in the project tag list; declared so the reference resolves inside the assembled package, while every step the counted argument uses remains inside the retained spans.\par\medskip
\ifdefined\hypertarget\hypertarget{sourcepointer-21}{}\fi
\noindent\textit{Source reference.} Original lemma 021I: sources/stacks-chow/tags.txt, source lines 2671--2671. Cross-chapter pointer whose target carries the canonical tag 021I in the project tag list; declared so the reference resolves inside the assembled package, while every step the counted argument uses remains inside the retained spans.\par\medskip
\ifdefined\hypertarget\hypertarget{sourcepointer-22}{}\fi
\noindent\textit{Source reference.} Original lemma 03F3: sources/stacks-chow/tags.txt, source lines 4478--4478. Cross-chapter pointer whose target carries the canonical tag 03F3 in the project tag list; declared so the reference resolves inside the assembled package, while every step the counted argument uses remains inside the retained spans.\par\medskip
\ifdefined\hypertarget\hypertarget{sourcepointer-23}{}\fi
\noindent\textit{Source reference.} Original lemma 08FI: sources/stacks-chow/tags.txt, source lines 10711--10711. Cross-chapter pointer whose target carries the canonical tag 08FI in the project tag list; declared so the reference resolves inside the assembled package, while every step the counted argument uses remains inside the retained spans.\par\medskip
\ifdefined\hypertarget\hypertarget{sourcepointer-24}{}\fi
\noindent\textit{Source reference.} Original lemma 0758: sources/stacks-chow/tags.txt, source lines 9122--9122. Cross-chapter pointer whose target carries the canonical tag 0758 in the project tag list; declared so the reference resolves inside the assembled package, while every step the counted argument uses remains inside the retained spans.\par\medskip
\ifdefined\hypertarget\hypertarget{sourcepointer-25}{}\fi
\noindent\textit{Source reference.} Original lemma 095S: sources/stacks-chow/tags.txt, source lines 11595--11595. Cross-chapter pointer whose target carries the canonical tag 095S in the project tag list; declared so the reference resolves inside the assembled package, while every step the counted argument uses remains inside the retained spans.\par\medskip
\ifdefined\hypertarget\hypertarget{sourcepointer-26}{}\fi
\noindent\textit{Source reference.} Original lemma 07U6: sources/stacks-chow/tags.txt, source lines 9961--9961. Cross-chapter pointer whose target carries the canonical tag 07U6 in the project tag list; declared so the reference resolves inside the assembled package, while every step the counted argument uses remains inside the retained spans.\par\medskip
\ifdefined\hypertarget\hypertarget{sourcepointer-27}{}\fi
\noindent\textit{Source reference.} Original lemma 015E: sources/stacks-chow/tags.txt, source lines 1505--1505. Cross-chapter pointer whose target carries the canonical tag 015E in the project tag list; declared so the reference resolves inside the assembled package, while every step the counted argument uses remains inside the retained spans.\par\medskip
\begin{lemma}
\label{lemma-describe-pullback}
Let $S$ be a scheme. Let $X$ be an algebraic space over $S$.
Let $\mathcal{F}$ be a sheaf on $X_\etale$. Then
$\pi_X^{-1}\mathcal{F}$ is given by the rule
$$
(\pi_X^{-1}\mathcal{F})(Y) = \Gamma(Y_\etale, f_{small}^{-1}\mathcal{F})
$$
for $f : Y \to X$ in $(\textit{Spaces}/X)_\etale$.
Moreover, $\pi_Y^{-1}\mathcal{F}$ satisfies the
sheaf condition with respect to smooth, syntomic, fppf, fpqc, and ph coverings.
\end{lemma}

\begin{proof}
Since pullback is transitive and $f_{small} = \pi_X \circ i_f$
(see above) we see that
$i_f^{-1} \pi_X^{-1}\mathcal{F} = f_{small}^{-1}\mathcal{F}$.
This shows that $\pi_X^{-1}$ has the description given in the lemma.

\medskip\noindent
To prove that $\pi_X^{-1}\mathcal{F}$ is a sheaf for the ph topology
it suffices by Topologies on Spaces, Lemma
\ref{spaces-topologies-lemma-characterize-sheaf}
to show that for a surjective proper morphism
$V \to U$ of algebraic spaces over $X$ we have
$(\pi_X^{-1}\mathcal{F})(U)$ is the equalizer of the two maps
$(\pi_X^{-1}\mathcal{F})(V) \to (\pi_X^{-1}\mathcal{F})(V \times_U V)$.
This we have seen in Lemma \ref{lemma-surjective-proper}.

\medskip\noindent
The case of smooth, syntomic, fppf coverings follows from the case
of ph coverings by Topologies on Spaces, Lemma
\ref{spaces-topologies-lemma-zariski-etale-smooth-syntomic-fppf-ph}.

\medskip\noindent
Let $\mathcal{U} = \{U_i \to U\}_{i \in I}$ be an fpqc covering of algebraic
spaces over $X$. Let $s_i \in (\pi_X^{-1}\mathcal{F})(U_i)$ be sections
which agree over $U_i \times_U U_j$. We have to prove there exists a unique
$s \in (\pi_X^{-1}\mathcal{F})(U)$ restricting to $s_i$ over $U_i$.
Case I: $U$ and $U_i$ are schemes. This case follows from
\'Etale Cohomology, Lemma \ref{etale-cohomology-lemma-describe-pullback}.
Case II: $U$ is a scheme. Here we choose surjective \'etale morphisms
$T_i \to U_i$ where $T_i$ is a scheme. Then $\mathcal{T} = \{T_i \to U\}$ is an
fpqc covering by schemes and by case I the result holds for $\mathcal{T}$.
We omit the verification that this implies the result for $\mathcal{U}$.
Case III: general case. Let $W \to U$ be a surjective \'etale
morphism, where $W$ is a scheme. Then $\mathcal{W} = \{U_i \times_U W \to W\}$
is an fpqc covering (by algebraic spaces) of the scheme $W$.
By case II the result hold for $\mathcal{W}$.
We omit the verification that this implies the result for $\mathcal{U}$.
\end{proof}

\begin{lemma}
\label{lemma-proper-base-change-f-star}
Let $S$ be a scheme. Let $f : X \to Y$ and $g : Y' \to Y$
be a morphisms of algebraic spaces over $S$. Assume $f$ is proper.
Set $X' = Y' \times_Y X$ with projections $f' : X' \to Y'$ and $g' : X' \to X$.
Let $\mathcal{F}$ be any sheaf on $X_\etale$. Then
$g^{-1}f_*\mathcal{F} = f'_*(g')^{-1}\mathcal{F}$.
\end{lemma}

\begin{proof}
The question is \'etale local on $Y'$. Choose a scheme $V$ and a surjective
\'etale morphism $V \to Y$. Choose a scheme $V'$ and a surjective \'etale
morphism $V' \to V \times_Y Y'$. Then we may replace $Y'$ by $V'$ and
$Y$ by $V$. Hence we may assume $Y$ and $Y'$ are schemes.
Then we may work Zariski locally on $Y$ and $Y'$ and hence we may
assume $Y$ and $Y'$ are affine schemes.

\medskip\noindent
Assume $Y$ and $Y'$ are affine schemes. Choose a surjective proper morphism
$h_1 : X_1 \to X$ where $X_1$ is a scheme, see
Cohomology of Spaces, Lemma \ref{spaces-cohomology-lemma-weak-chow}.
Set $X_2 = X_1 \times_X X_1$ and denote
$h_2 : X_2 \to X$ the structure morphism. Observe this is a scheme.
By the case of schemes
(\'Etale Cohomology, Lemma
\ref{etale-cohomology-lemma-proper-base-change-f-star})
we know the lemma is true for the cartesian diagrams
$$
\vcenter{
\xymatrix{
X'_1 \ar[r] \ar[d] & X_1 \ar[d] \\
Y' \ar[r] & Y
}
}
\quad\text{and}\quad
\vcenter{
\xymatrix{
X'_2 \ar[r] \ar[d] & X_2 \ar[d] \\
Y' \ar[r] & Y
}
}
$$
and the sheaves $\mathcal{F}_i = (X_i \to X)^{-1}\mathcal{F}$.
By Lemma \ref{lemma-surjective-proper} we have an exact sequence
$0 \to \mathcal{F} \to h_{1, *}\mathcal{F}_1 \to h_{2, *}\mathcal{F}_2$
and similarly for $(g')^{-1}\mathcal{F}$ because
$X'_2 = X'_1 \times_{X'} X'_1$. Hence we conlude that the
lemma is true (some details omitted).
\end{proof}

\begin{lemma}
\label{lemma-compare-higher-direct-image}
Let $S$ be a scheme.
Let $f : Y \to X$ be a morphism of algebraic spaces over $S$.
\begin{enumerate}
\item For $K$ in $D((\textit{Spaces}/Y)_\etale)$ we have
$
(Rf_{big, *}K)|_{X_\etale} = Rf_{small, *}(K|_{Y_\etale})
$
in $D(X_\etale)$.
\item For $K$ in $D((\textit{Spaces}/Y)_\etale, \mathcal{O})$ we have
$
(Rf_{big, *}K)|_{X_\etale} = Rf_{small, *}(K|_{Y_\etale})
$
in $D(\textit{Mod}(X_\etale, \mathcal{O}_X))$.
\end{enumerate}
More generally, let $g : X' \to X$ be an object of
$(\textit{Spaces}/X)_\etale$. Consider the fibre product
$$
\xymatrix{
Y' \ar[r]_{g'} \ar[d]_{f'} & Y \ar[d]^f \\
X' \ar[r]^g & X
}
$$
Then
\begin{enumerate}
\item[(3)] For $K$ in $D((\textit{Spaces}/Y)_\etale)$ we have
$i_g^{-1}(Rf_{big, *}K) = Rf'_{small, *}(i_{g'}^{-1}K)$
in $D(X'_\etale)$.
\item[(4)] For $K$ in $D((\textit{Spaces}/Y)_\etale, \mathcal{O})$ we have
$i_g^*(Rf_{big, *}K) = Rf'_{small, *}(i_{g'}^*K)$
in $D(\textit{Mod}(X'_\etale, \mathcal{O}_{X'}))$.
\item[(5)] For $K$ in $D((\textit{Spaces}/Y)_\etale)$ we have
$g_{big}^{-1}(Rf_{big, *}K) = Rf'_{big, *}((g'_{big})^{-1}K)$
in $D((\textit{Spaces}/X')_\etale)$.
\item[(6)] For $K$ in $D((\textit{Spaces}/Y)_\etale, \mathcal{O})$ we have
$g_{big}^*(Rf_{big, *}K) = Rf'_{big, *}((g'_{big})^*K)$
in $D(\textit{Mod}(X'_\etale, \mathcal{O}_{X'}))$.
\end{enumerate}
\end{lemma}

\begin{proof}
Part (1) follows from
Lemma \ref{lemma-compare-injectives}
and (\ref{equation-compare-big-small})
on choosing a K-injective complex of abelian sheaves representing $K$.

\medskip\noindent
Part (3) follows from Lemma \ref{lemma-compare-injectives}
and Topologies, Lemma
\ref{topologies-lemma-morphism-big-small-cartesian-diagram-etale}
on choosing a K-injective complex of abelian sheaves representing $K$.

\medskip\noindent
Part (5) is Cohomology on Sites, Lemma
\ref{sites-cohomology-lemma-localize-cartesian-square}.

\medskip\noindent
Part (6) is  Cohomology on Sites, Lemma
\ref{sites-cohomology-lemma-localize-cartesian-square-modules}.

\medskip\noindent
Part (2) can be proved as follows. Above we have seen
that $\pi_X \circ f_{big} = f_{small} \circ \pi_Y$ as morphisms
of ringed sites. Hence we obtain
$R\pi_{X, *} \circ Rf_{big, *} = Rf_{small, *} \circ R\pi_{Y, *}$
by Cohomology on Sites, Lemma
\ref{sites-cohomology-lemma-derived-pushforward-composition}.
Since the restriction functors $\pi_{X, *}$ and $\pi_{Y, *}$
are exact, we conclude.

\medskip\noindent
Part (4) follows from part (6) and part (2) applied to $f' : Y' \to X'$.
\end{proof}

\begin{lemma}
\label{lemma-proper-base-change}
Let $S$ be a scheme. Let
$$
\xymatrix{
X' \ar[r]_{g'} \ar[d]_{f'} & X \ar[d]^f \\
Y' \ar[r]^g & Y
}
$$
be a cartesian square of algebraic spaces over $S$. Assume $f$ is proper.
Let $E \in D^+(X_\etale)$ have torsion cohomology sheaves.
Then the base change map $g^{-1}Rf_*E \to Rf'_*(g')^{-1}E$
is an isomorphism.
\end{lemma}

\begin{proof}
This is a simple consequence of the proper base change theorem
(Theorem \ref{theorem-proper-base-change}) using the spectral
sequences
$$
E_2^{p, q} = R^pf_*H^q(E)
\quad\text{and}\quad
{E'}_2^{p, q} = R^pf'_*(g')^{-1}H^q(E)
$$
converging to $R^nf_*E$ and $R^nf'_*(g')^{-1}E$.
The spectral sequences are constructed in
Derived Categories, Lemma \ref{derived-lemma-two-ss-complex-functor}.
Some details omitted.
\end{proof}

\begin{lemma}
\label{lemma-surjective-proper}
Let $S$ be a scheme. Let $f : Y \to X$ be a surjective proper morphism
of algebraic spaces over $S$. Let $\mathcal{F}$ be a sheaf on $X_\etale$.
Then $\mathcal{F} \to f_*f^{-1}\mathcal{F}$ is injective with
image the equalizer of the two maps
$f_*f^{-1}\mathcal{F} \to g_*g^{-1}\mathcal{F}$ where
$g$ is the structure morphism $g : Y \times_X Y \to X$.
\end{lemma}

\begin{proof}
For any surjective morphism $f : Y \to X$ of algebraic spaces over $S$,
the map $\mathcal{F} \to f_*f^{-1}\mathcal{F}$ is injective.
Namely, if $\overline{x}$ is a geometric point of $X$, then we
choose a geometric point $\overline{y}$ of $Y$ lying over $\overline{x}$
and we consider
$$
\mathcal{F}_{\overline{x}} \to
(f_*f^{-1}\mathcal{F})_{\overline{x}} \to
(f^{-1}\mathcal{F})_{\overline{y}} = \mathcal{F}_{\overline{x}}
$$
See Properties of Spaces, Lemma \ref{spaces-properties-lemma-stalk-pullback}
for the last equality.

\medskip\noindent
The second statement is local on $X$ in the \'etale topology, hence we may
and do assume $Y$ is an affine scheme.

\medskip\noindent
Choose a surjective proper morphism $Z \to Y$ where $Z$ is a scheme, see
Cohomology of Spaces, Lemma \ref{spaces-cohomology-lemma-weak-chow}.
The result for $Z \to X$ implies the result for $Y \to X$.
Since $Z \to X$ is a surjective proper morphism of schemes
and hence a ph covering
(Topologies, Lemma \ref{topologies-lemma-surjective-proper-ph})
the result for $Z \to X$ follows from
\'Etale Cohomology, Lemma \ref{etale-cohomology-lemma-describe-pullback-pi-ph}
(in fact it is in some sense equivalent to this lemma).
\end{proof}

\begin{equation}
\label{equation-compare-big-small}
(f_{big, *}\mathcal{F})|_{X_\etale} =
f_{small, *}(\mathcal{F}|_{Y_\etale})
\end{equation}

\begin{lemma}
\label{lemma-compare-injectives}
Let $S$ be a scheme.
Let $Y \to X$ be a morphism of $(\textit{Spaces}/S)_\etale$.
\begin{enumerate}
\item If $\mathcal{I}$ is injective in
$\textit{Ab}((\textit{Spaces}/X)_\etale)$, then
\begin{enumerate}
\item $i_f^{-1}\mathcal{I}$ is injective in $\textit{Ab}(Y_\etale)$,
\item $\mathcal{I}|_{X_\etale}$ is injective in $\textit{Ab}(X_\etale)$,
\end{enumerate}
\item If $\mathcal{I}^\bullet$ is a K-injective complex
in $\textit{Ab}((\textit{Spaces}/X)_\etale)$, then
\begin{enumerate}
\item $i_f^{-1}\mathcal{I}^\bullet$ is a K-injective complex in
$\textit{Ab}(Y_\etale)$,
\item $\mathcal{I}^\bullet|_{X_\etale}$ is a K-injective complex in
$\textit{Ab}(X_\etale)$,
\end{enumerate}
\end{enumerate}
The corresponding statements for modules do not hold.
\end{lemma}

\begin{proof}
Parts (1)(b) and (2)(b)
follow formally from the fact that the restriction functor
$\pi_{X, *} = i_X^{-1}$ is a right adjoint of the exact functor
$\pi_X^{-1}$, see
Homology, Lemma \ref{homology-lemma-adjoint-preserve-injectives} and
Derived Categories, Lemma \ref{derived-lemma-adjoint-preserve-K-injectives}.

\medskip\noindent
Parts (1)(a) and (2)(a) can be seen in two ways. First proof: We can use
that $i_f^{-1}$ is a right adjoint of the exact functor $i_{f, !}$.
This functor is constructed in
Topologies, Lemma \ref{topologies-lemma-put-in-T-etale}
for sheaves of sets and for abelian sheaves in
Modules on Sites, Lemma \ref{sites-modules-lemma-g-shriek-adjoint}.
It is shown in Modules on Sites, Lemma
\ref{sites-modules-lemma-exactness-lower-shriek} that it is exact.
Second proof. We can use that $i_f = i_Y \circ f_{big}$ as is shown
in Topologies, Lemma \ref{topologies-lemma-morphism-big-small-etale}.
Since $f_{big}$ is a localization, we see that pullback by it
preserves injectives and K-injectives, see
Cohomology on Sites, Lemmas \ref{sites-cohomology-lemma-cohomology-of-open} and
\ref{sites-cohomology-lemma-restrict-K-injective-to-open}.
Then we apply the already proved parts (1)(b) and (2)(b)
to the functor $i_Y^{-1}$ to conclude.

\medskip\noindent
To see a counter example for the case of modules we refer to
\'Etale Cohomology, Lemma \ref{etale-cohomology-lemma-compare-injectives}.
\end{proof}

\begin{theorem}
\label{theorem-proper-base-change}
Let $S$ be a scheme. Let
$$
\xymatrix{
X' \ar[r]_{g'} \ar[d]_{f'} & X \ar[d]^f \\
Y' \ar[r]^g & Y
}
$$
be a cartesian square of algebraic spaces over $S$.
Assume $f$ is proper.
Let $\mathcal{F}$ be an abelian torsion sheaf on $X_\etale$.
Then the base change map
$$
g^{-1}Rf_*\mathcal{F} \longrightarrow Rf'_*(g')^{-1}\mathcal{F}
$$
is an isomorphism.
\end{theorem}

\begin{proof}
This proof repeats a few of the arguments given in the proof of the
proper base change theorem for schemes. See
\'Etale Cohomology, Section \ref{etale-cohomology-section-proper-base-change}
for more details.

\medskip\noindent
The statement is \'etale local on $Y'$ and $Y$, hence we may assume
both $Y$ and $Y'$ are affine schemes. Observe that this in particular
proves the theorem in case $f$ is representable (we will use this
below).

\medskip\noindent
For every $n \geq 1$ let $\mathcal{F}[n]$ be the subsheaf of sections
of $\mathcal{F}$ annihilated by $n$. Then
$\mathcal{F} = \colim \mathcal{F}[n]$. By 
Cohomology of Spaces, Lemma \ref{spaces-cohomology-lemma-colimit-cohomology}
the functors $g^{-1}R^pf_*$ and $R^pf'_*(g')^{-1}$ commute
with filtered colimits. Hence it suffices to prove the theorem
if $\mathcal{F}$ is killed by $n$.

\medskip\noindent
Let $\mathcal{F} \to \mathcal{I}^\bullet$ be a resolution by
injective sheaves of $\mathbf{Z}/n\mathbf{Z}$-modules.
Observe that
$g^{-1}f_*\mathcal{I}^\bullet = f'_*(g')^{-1}\mathcal{I}^\bullet$
by Lemma \ref{lemma-proper-base-change-f-star}.
Applying Leray's acyclicity lemma
(Derived Categories, Lemma \ref{derived-lemma-leray-acyclicity})
we conclude it suffices to prove
$R^pf'_*(g')^{-1}\mathcal{I}^m = 0$ for $p > 0$ and $m \in \mathbf{Z}$.

\medskip\noindent
Choose a surjective proper morphism
$h : Z \to X$ where $Z$ is a scheme, see
Cohomology of Spaces, Lemma \ref{spaces-cohomology-lemma-weak-chow}.
Choose an injective map $h^{-1}\mathcal{I}^m \to \mathcal{J}$
where $\mathcal{J}$ is an injective sheaf of
$\mathbf{Z}/n\mathbf{Z}$-modules on $Z_\etale$.
Since $h$ is surjective the map $\mathcal{I}^m \to h_*\mathcal{J}$
is injective (see Lemma \ref{lemma-surjective-proper}).
Since $\mathcal{I}^m$ is injective we see that $\mathcal{I}^m$
is a direct summand of $h_*\mathcal{J}$. Thus it suffices
to prove the desired vanishing for $h_*\mathcal{J}$.

\medskip\noindent
Denote $h'$ the base change by $g$ and denote $g'' : Z' \to Z$
the projection. There is a spectral sequence
$$
E_2^{p, q} = R^pf'_* R^qh'_* (g'')^{-1}\mathcal{J}
$$
converging to $R^{p + q}(f' \circ h')_*(g'')^{-1}\mathcal{J}$.
Since $h$ and $f \circ h$ are representable (by schemes)
we know the result we want holds for them. Thus in the
spectral sequence we see that $E_2^{p, q} = 0$ for $q > 0$
and $R^{p + q}(f' \circ h')_*(g'')^{-1}\mathcal{J} = 0$
for $p + q > 0$. It follows that $E_2^{p, 0} = 0$ for $p > 0$.
Now
$$
E_2^{p, 0} = R^pf'_* h'_* (g'')^{-1}\mathcal{J} =
R^pf'_* (g')^{-1}h_*\mathcal{J}
$$
by Lemma \ref{lemma-proper-base-change-f-star}. This finishes the proof.
\end{proof}

\begin{lemma}
\label{lemma-compare-higher-direct-image-proper}
Let $S$ be a scheme.
Let $f : Y \to X$ be a proper morphism of algebraic spaces over $S$.
Then we have
\begin{enumerate}
\item $\pi_X^{-1} \circ f_{small, *} = f_{big, *} \circ \pi_Y^{-1}$
as functors $\Sh(Y_\etale) \to \Sh((\textit{Spaces}/X)_\etale)$,
\item $\pi_X^{-1}Rf_{small, *}K = Rf_{big, *}\pi_Y^{-1}K$
for $K$ in $D^+(Y_\etale)$ whose cohomology sheaves are torsion, and
\item $\pi_X^{-1}Rf_{small, *}K = Rf_{big, *}\pi_Y^{-1}K$
for all $K$ in $D(Y_\etale)$ if $f$ is finite.
\end{enumerate}
\end{lemma}

\begin{proof}
Proof of (1). Let $\mathcal{F}$ be a sheaf on $Y_\etale$.
Let $g : X' \to X$ be an object of $(\textit{Spaces}/X)_\etale$.
Consider the fibre product
$$
\xymatrix{
Y' \ar[r]_{f'} \ar[d]_{g'} & X' \ar[d]^g \\
Y \ar[r]^f & X
}
$$
Then we have
$$
(f_{big, *}\pi_Y^{-1}\mathcal{F})(X') =
(\pi_Y^{-1}\mathcal{F})(Y') =
((g'_{small})^{-1}\mathcal{F})(Y')  =
(f'_{small, *}(g'_{small})^{-1}\mathcal{F})(X')
$$
the second equality by Lemma \ref{lemma-describe-pullback}.
On the other hand
$$
(\pi_X^{-1}f_{small, *}\mathcal{F})(X') =
(g_{small}^{-1}f_{small, *}\mathcal{F})(X')
$$
again by Lemma \ref{lemma-describe-pullback}.
Hence by proper base change for sheaves of sets
(Lemma \ref{lemma-proper-base-change-f-star})
we conclude the two sets are canonically isomorphic.
The isomorphism is compatible with restriction mappings
and defines an isomorphism
$\pi_X^{-1}f_{small, *}\mathcal{F} = f_{big, *}\pi_Y^{-1}\mathcal{F}$.
Thus an isomorphism of functors
$\pi_X^{-1} \circ f_{small, *} = f_{big, *} \circ \pi_Y^{-1}$.

\medskip\noindent
Proof of (2). There is a canonical base change map
$\pi_X^{-1}Rf_{small, *}K \to Rf_{big, *}\pi_Y^{-1}K$
for any $K$ in $D(Y_\etale)$, see
Cohomology on Sites, Remark \ref{sites-cohomology-remark-base-change}.
To prove it is an isomorphism, it suffices to prove the pull back of
the base change map by $i_g : \Sh(X'_\etale) \to \Sh((\Sch/X)_\etale)$
is an isomorphism for any object $g : X' \to X$ of $(\Sch/X)_\etale$.
Let $T', g', f'$ be as in the previous paragraph.
The pullback of the base change map is
\begin{align*}
g_{small}^{-1}Rf_{small, *}K
& =
i_g^{-1}\pi_X^{-1}Rf_{small, *}K \\
& \to
i_g^{-1}Rf_{big, *}\pi_Y^{-1}K \\
& =
Rf'_{small, *}(i_{g'}^{-1}\pi_Y^{-1}K) \\
& =
Rf'_{small, *}((g'_{small})^{-1}K)
\end{align*}
where we have used $\pi_X \circ i_g = g_{small}$,
$\pi_Y \circ i_{g'} = g'_{small}$, and
Lemma \ref{lemma-compare-higher-direct-image}.
This map is an isomorphism by the proper base change theorem
(Lemma \ref{lemma-proper-base-change}) provided $K$ is bounded
below and the cohomology sheaves of $K$ are torsion.

\medskip\noindent
Proof of (3). If $f$ is finite, then the functors
$f_{small, *}$ and $f_{big, *}$ are exact. This follows
from Cohomology of Spaces, Lemma
\ref{spaces-cohomology-lemma-finite-higher-direct-image-zero}
for $f_{small}$. Since any base change $f'$ of $f$ is finite too,
we conclude from Lemma \ref{lemma-compare-higher-direct-image} part (3)
that $f_{big, *}$ is exact too (as the higher derived functors are zero).
Thus this case follows from part (1).
\end{proof}

\end{document}
